\clearpage\begingroup\SourceLeaf{387b}\ArithmeticHead{56}{L}{STÄCKEL, GAUSS ALS GEOMETER.}
\renewcommand{\thefootnote}{\arabic{footnote})}\setcounter{footnote}{1}
\begin{center}
\textls[100]{Abschnitt III.}\par\medskip
\textbf{Die komplexen Grössen in ihrer Beziehung zur Geometrie.}\par\medskip
\textbf{14.}\par\smallskip
\textls[100]{Kreisteilung.}
\end{center}
\noindent\hspace*{\parindent}Die »Darstellung der imaginären Grössen in den Relationen der Punkte
in plano« (Brief an \textsc{Drobisch} vom 14. August 1834, W. X\textsubscript{1}, S. 106) hat nicht
nur für die arithmetischen und funktionentheoretischen, sondern auch für die
geometrischen Untersuchungen von \textsc{Gauss} eine so grosse Bedeutung, dass den
Beziehungen der komplexen Grössen zur Raumlehre ein besonderer Abschnitt
dieses Aufsatzes gewidmet werden soll; in ihm sollen die Ausführungen, die
in den Aufsätzen über \textsc{Gauss}’ Arbeiten zur Zahlentheorie, Funktionentheorie
und Algebra enthalten sind, wieder aufgenommen und ergänzt werden.
\par
Schon sehr früh hat \textsc{Gauss} versucht, um einen von ihm gern gebrauchten
Ausdruck anzuwenden, in die \textls[100]{Metaphysik} der imaginären Grössen einzudringen.
In der Selbstanzeige der zweiten Abhandlung \textit{über die biquadratischen Reste}
vom Jahre 1831 sagt er, dass er »diesen hochwichtigen Teil der Mathematik
seit vielen Jahren betrachtet habe« (W. II, S. 175), und in dem Briefe an
\textsc{Drobisch} vom 14. August 1834 freut er sich, dass dieser »auf seine schon fast
\clearpage\SourceLeaf{388}\ArithmeticHead{57}{R}{DIE KOMPLEXEN GRÖSSEN IN IHRER BEZIEHUNG ZUR GEOMETRIE.}\setcounter{footnote}{0}
\noindent seit 40 Jahren gehegten Grundansichten über die imaginären Grössen eingegangen
sei« (W. X\textsubscript{1}, S. 106). Als solche Grundansichten wird man wohl
erstens die Erkenntnis zu bezeichnen haben, dass »den komplexen Grössen
das völlig gleiche Bürgerrecht mit den reellen Grössen eingeräumt werden
müsse« (W. II, S. 171), und zweitens, dass diese Grössen »ebenso gut wie die
negativen ihre reale gegenständliche Bedeutung haben« (W. X\textsubscript{1}, S. 405), die
sich in ihrer »Versinnlichung durch die Punkte einer unbegrenzten Ebene«
(W. X\textsubscript{1}, S. 407) kund gibt.
\par
Wird man durch die vorstehenden Angaben von \textsc{Gauss} etwa auf die
Jahre 1795 und 1796 zurückgeführt, so kann als Bestätigung eine Stelle der
\textit{Disquisitiones arithmeticae} dienen, und zwar aus dem dritten Abschnitt, der nach
\textsc{Bachmann} (W. X\textsubscript{2}, Abh. 1, S. 6) im Wesentlichen bereits 1796 entstanden und
1797 niedergeschrieben worden ist (der Druck der \textit{Disquisitiones} begann im
April 1798 und hat mit verschiedenen Unterbrechungen bis 1801 gedauert).
Dort sagt \textsc{Gauss}, er wolle auf die Lehre von den imaginären Indizes, zu denen
man bei Moduln ohne primitive Wurzeln seine Zuflucht nehmen muss, bei
einer anderen Gelegenheit eingehen, »wenn wir es vielleicht unternehmen
werden, die Lehre von den imaginären Grössen, die wenigstens nach unserem
Urteil bis jetzt von Niemandem auf klare Begriffe zurückgeführt ist, ausführlicher
zu behandeln« (W I, S. 71).
\par
Im \textit{Tagebuch}, das mit dem März 1796 beginnt, findet sich keine Aufzeichnung,
die man mit einer solchen Absicht in Verbindung bringen könnte.
Wohl aber zeigt gerade die erste Eintragung, dass \textsc{Gauss} in der vorhergehenden
Zeit mit imaginären Grössen zu tun gehabt hatte, denn er verkündet
hier, dass er die geometrische Siebzehnteilung des Kreisumfanges entdeckt
habe, das heisst, wie wir aus dem Briefe an \textsc{Gerling} vom 6. Januar
1819 (W. X\textsubscript{1}, S. 125) wissen, die Auflösung der zugehörigen Kreisteilungsgleichung
mittels wiederholter Ausziehung von Quadratwurzeln, und zwar hatte
\textsc{Gauss}, nach den Angaben in demselben Briefe, schon während seines ersten
Semesters in Göttingen, das Oktober 1795 begann, die Kreisteilungsgleichungen
für einen beliebigen Primzahlgrad untersucht.
\par
Dass die Teilung des Kreisumfanges in \(n\) gleiche Stücke mittels imaginärer
Grössen auf die Lösung der Gleichung \(x^n-1=0\) zurückgeführt werden
\par\nobreak{\small X\textsubscript{2} Abh. 4.\hfill 8}
\clearpage\SourceLeaf{389}\ArithmeticHead{58}{L}{STÄCKEL, GAUSS ALS GEOMETER.}\setcounter{footnote}{0}
\noindent kann, ist eine Einsicht, die man \textsc{Cotes}\footnote{R. \textsc{Cotes}, \textit{Harmonia mensurarum, sive analysis et synthesis per rationum et angulorum mensuras promota}, Cambridge 1722.} und \textsc{Moivre}\footnote{A. \textsc{de Moivre}, \textit{Miscellanea analytica}, London 1730.} verdankt, die aber
erst durch \textsc{Euler} geklärt und sichergestellt worden ist\footnote{L. \textsc{Euler}, \textit{Introductio in analysin}, Lausanne und Genf 1748, siehe besonders t. I, cap. 8: \textit{De quantitatibus transcendentibus ex circulo ortis}.}. Später hat sich
\textsc{Vandermonde} mit der Auflösung solcher Gleichungen mittels Wurzelziehens
befasst. Es ist sehr wahrscheinlich, dass \textsc{Gauss} bei der Abfassung der \textit{Disquisitiones
arithmeticae} dessen Abhandlung gekannt hat, denn in dem Briefe an
\textsc{Olbers} vom 12. Oktober 1802 sagt er, dass wir über die Geometria situs »nur
einige Fragmente von \textsc{Euler} und einem von mir sehr hochgeschätzten Geometer
\textsc{Vandermonde} haben« (Br. G.-O., 1, S. 103). Die Abhandlung über
Geometria situs steht aber in demselben Bande der Pariser Denkschriften für
das Jahr 1771 wie die Abhandlung über die Auflösung der algebraischen
Gleichungen\footnote{\textsc{Ch. A. Vandermonde}, \textit{Remarques sur les problèmes de situation}, Histoire de l’Acad., année 1771, Paris 1774, Mémoires, S. 566; \textit{Sur la résolution des équations}; ebenda, S. 365; die letztere Abhandlung ist in deutscher Sprache herausgegeben von C. \textsc{Itzigsohn}, \textsc{Vandermonde}, \textit{Abhandlungen aus der reinen Mathematik}, Berlin 1887. Auf S. 375 behauptet \textsc{Vandermonde}, die Gleichung \(x^n-1=0\) sei für jeden Grad \(n\) durch Wurzelziehen lösbar und führt die Rechnungen für einige Fälle durch, im Besonderen für \(n=11\). Für die Exponenten \(n\leqq10\) hatte schon \textsc{Euler}, \textit{De extractione radicum ex quantitatibus irrationalibus}, Comment. acad. sc. Petrop. 13 (1741/3) 1751, § 39 bis 48, Opera omnia, ser. I, vol. 6, S. 31, die Wurzeln mittels blosser Wurzelziehungen dargestellt; dagegen, meint er, führe der Fall \(n=11\) auf eine Gleichung fünften Grades, deren Lösung noch verborgen sei.}.
\par
Nachdem \textsc{Gauss} im Art. 337 der \textit{Disquisitiones arithmeticae} (W. I, S. 414)
bemerkt hat, die trigonometrischen Funktionen der Bögen
\[
2k\pi/n\qquad(k=0,1,2,\mathord{.}\mathord{.}\mathord{.},n-1)
\]
\noindent seien die Wurzeln von Gleichungen \(n\)-ten Grades, fährt er fort: »Jedoch ist
keine dieser Gleichungen so leicht zu behandeln und für unseren Zweck so
geeignet, wie diese: \(x^n-1=0\), deren Wurzeln bekanntlich mit den Wurzeln
jener aufs engste verbunden sind. Wenn man nämlich der Kürze halber \(i\)
für die imaginäre Grösse \(\surd{-1}\) schreibt, so werden die Wurzeln der Gleichung
\(x^n-1=0\) durch
\[
\cos 2k\pi/n+i\sin 2k\pi/n
\]
\noindent dargestellt, wo für \(k\) alle Zahlen \(0,1,2,\mathord{.}\mathord{.}\mathord{.},n-1\) zu nehmen sind«.
\clearpage\SourceLeaf{390}\ArithmeticHead{59}{R}{DIE KOMPLEXEN GRÖSSEN IN IHRER BEZIEHUNG ZUR GEOMETRIE.}\setcounter{footnote}{0}
\noindent\hspace*{\parindent}Auf diese Art werden den Eckpunkten des regelmässigen \(n\)-Ecks, das
dem Kreise vom Halbmesser Eins eingeschrieben ist, die soeben angegebenen
komplexen Grössen zugeordnet. Die dabei auftretenden Grössen \(\cos 2k\pi/n\)
und \(\sin 2k\pi/n\) sind die rechtwinkligen kartesischen Koordinaten der betreffenden
Eckpunkte, wenn der Mittelpunkt des Kreises zum Anfangspunkt
gewählt und die Abszissenachse durch den Eckpunkt gelegt wird, für den
\(k=0\) ist. Mithin gelangt man in diesem Falle ganz unmittelbar zu der
\textsc{Gauss}schen Versinnlichung der komplexen Grössen durch die Punkte einer
Ebene.
\par
Dass die Betrachtung der Eckpunkte des \(n\)-Ecks \textsc{Gauss} geläufig war,
zeigt auch die Ausdrucksweise, ganze Zahlen seien »kongruent modulo \(n\)«,
wenn sie sich um Vielfache einer ganzen Zahl \(n\) unterscheiden: beim Durchlaufen
des Kreisumfangs entsprechen nämlich den Werten von \(k\), die mod. \(n\)
kongruent sind, dieselben Eckpunkte des \(n\)-Ecks, und so hat die Bezeichnung
»kongruent« ihre gute geometrische Bedeutung.
\par
Ob die geometrische Versinnlichung der komplexen Grössen den Untersuchungen
über die Kreisteilung entsprungen ist, lässt sich freilich nicht mit
Sicherheit entscheiden. Man könnte dagegen einwenden, dass auch bei \textsc{Euler}
Grössen der Form \(\cos\varphi+i\sin\varphi\) an mehr als einer Stelle in einer Weise auftreten,
die ihre geometrische Bedeutung nahe zu legen scheint, ohne dass es
dazu gekommen ist, und die Hauptsache liegt in dem Entschluss, die imaginären
Grössen als den reellen gleichberechtigt anzuerkennen. Vielleicht hat
\textsc{Gauss} diese Anerkennung durch die bereits erwähnte Einführung des Zeichens
\(i\) im art. 337 der \textit{Disq. arith.} andeuten wollen\footnote{Das Zeichen \(i\) für \(\surd{-1}\) findet sich \textls[100]{gelegentlich} schon bei \textsc{Euler}, nämlich in der am 5. Mai 1777 der Petersburger Akademie vorgelegten Abhandlung: \textit{De formulis differentialibus angularibus maxime irrationalibus, quas tamen per logarithmos et arcus circulares integrare licet}, die 1794 aus dem Nachlass im vierten Bande der \textit{Institutiones calculi integralis} abgedruckt ist, ed. tertia, Petersburg 1845, S. 184. \textsc{Gauss} hat das Zeichen \(i\) seit dem Jahre 1801 beständig angewandt und seinem Beispiel sind die Mathematiker gefolgt.}. Dass er sich in den Disquisitiones
wie in der Dissertation (1799) mit Andeutungen begnügte, ist wohl
teils aus seiner Scheu, strittige Dinge zu berühren, teils aus dem Umstande
zu erklären, dass er selbst, wenn auch seine »Grundansicht« feststand, die neue
Lehre noch nicht für reif hielt. In der Tat ist er erst nach einer langen
und harten Arbeit zu einer ihn befriedigenden Auffassung der imaginären
\par\nobreak\hfill{\small 8\textsuperscript{*}}
\clearpage\SourceLeaf{391a}\ArithmeticHead{60}{L}{STÄCKEL, GAUSS ALS GEOMETER.}\setcounter{footnote}{0}
\noindent Grössen gelangt. So schreibt er am 11. Dezember 1825 an \textsc{Hansen}, seine
Untersuchungen über die allgemeine Lehre von den krummen Flächen griffen
tief ein in die Metaphysik der Raumlehre, »und nur mit Mühe kann ich mich
von solchen daraus entspringenden Folgen, wie z. B. die wahre Metaphysik
der imaginären Grössen ist, losreissen. Der wahre Sinn des \(\surd{-1}\) steht mir
dabei mit grosser Lebendigkeit vor der Seele, aber es wird schwer sein, ihn
in Worte zu fassen, die immer nur ein vages, in der Luft schwebendes Bild
geben können« (Brief im \textsc{Gauss}-Archiv). In einer wahrscheinlich im Anschluss
an diesen Brief niedergeschriebenen Aufzeichnung \textit{Fragen zur Metaphysik der
Mathematik} (W. X\textsubscript{1}, S. 396) hat er versucht, seine Gedanken auszugestalten,
und man erkennt hier die Anfänge der Darstellung, die er in der Selbstanzeige
vom Jahre 1831 gegeben hat.
\clearpage\endgroup

